\section*{Capítulo \ref{ch:b2:living-soil} --- El suelo vivo}

\begin{solution}{exo:b2:living-soil:1}
O: hojarasca y materia orgánica parcialmente descompuesta, donde empieza la
descomposición. A: \omterm{def:b2:living-soil:soil}{suelo} mineral mezclado con \omterm{def:b2:living-soil:soil}{humus}, oscuro y de estructura
grumosa, que contiene la mayoría de las raíces, los organismos, los
nutrientes y el agua. B: el horizonte de acumulación, donde se depositan
la arcilla, los óxidos de hierro y los carbonatos lixiviados desde arriba.
C: material parental meteorizado, fragmentos de roca en una matriz fina,
donde se libera \omterm{def:b2:living-soil:soil}{suelo} mineral nuevo. R: lecho rocoso sin meteorizar.
\end{solution}

\begin{solution}{exo:b2:living-soil:2}
Las cargas negativas de las superficies de la arcilla y del \omterm{def:b2:living-soil:soil}{humus} retienen
cationes intercambiables que las raíces pueden tomar a cambio de protones.
Los granos de arena son de cuarzo, sin carga y gruesos, con poca superficie
por gramo y prácticamente sin \omterm{def:b2:living-soil:soil}{humus}. El encalado aporta carbonato de
calcio: el calcio desplaza los iones hidrógeno y aluminio que saturan los
sitios de intercambio de un \omterm{def:b2:living-soil:soil}{suelo} ácido, el carbonato neutraliza la acidez
y los sitios vuelven a llenarse con un catión nutritivo.
\end{solution}

\begin{solution}{exo:b2:living-soil:3}
Descomponedores: bacterias (\emph{Bacillus}) y hongos
(\emph{Trichoderma}); consumidores de microorganismos: \omterm{def:b2:unicellular-diversity:protist}{protistas} (amebas)
y nematodos bacterívoros; detritívoros: lombrices, milpiés, colémbolos,
cochinillas de la humedad; depredadores: ácaros, nematodos depredadores,
ciempiés, carábidos; y topos o musarañas en la cima. Los consumidores de
bacterias, con un $\mathrm{C:N}$ más alto que el de sus presas, excretan el
nitrógeno sobrante en forma de amonio: mineralizan el nitrógeno que las
bacterias habían retenido.
\end{solution}

\begin{solution}{exo:b2:living-soil:4}
Arbuscular: el hongo (un glomeromiceto) penetra en las células del córtex
de la raíz y forma arbúsculos; se encuentra en la mayoría de las hierbas y
los cultivos y en los árboles tropicales. Ectomicorrícica: el hongo
(basidiomicetos y ascomicetos, muchos de ellos con setas) envuelve el ápice
de la raíz y crece entre las células sin penetrar en ellas; se encuentra en
los árboles templados y boreales. En ambos casos, la planta da azúcar y el
hongo da fosfato, nitrógeno, agua, micronutrientes y cierta protección
frente a los patógenos.
\end{solution}

\begin{solution}{exo:b2:living-soil:5}
$L^{*} = I/k = 4/0.8 = \qty{5}{t/ha}$; \omterm{def:b2:biogeochemical-cycles:reservoir}{tiempo de residencia}, $1/k = 1.25$
años; pérdida del \qty{90}{\%} al cabo de $\ln 10/0.8 = 2.9$ años.
\end{solution}

\begin{solution}{exo:b2:living-soil:6}
Por cada \qty{100}{kg} de carbono comido, los microorganismos retienen
\qty{40}{kg} y necesitan $40/8 = \qty{5}{kg}$ de nitrógeno. La paja aporta
$100/80 =
\qty{1.25}{kg}$: inmovilización de \qty{3.75}{kg} por cada \qty{100}{kg}
de carbono. El trébol aporta $100/15 = \qty{6.7}{kg}$: mineralización
de \qty{1.7}{kg}.
\end{solution}

\begin{solution}{exo:b2:living-soil:7}
\qty{20}{t/ha} de tierra al año a través de las lombrices. Los
\qty{20}{cm} superiores pesan $0.2\times 10^{4}\times 1.3 = \qty{2600}{t/ha}$:
130 años para un paso. Las diez toneladas por acre de Darwin son
\qty{25}{t/ha}, del mismo orden --- su ``cada pocos años'' se refería al
mantillo superficial, de unos pocos centímetros, no a toda la capa
superficial.
\end{solution}

\begin{solution}{exo:b2:living-soil:8}
Bacterias: $10^{9}\times 10^{-12} = \qty{1}{mg}$ por gramo. Hifas:
volumen $\pi(2\times 10^{-4})^{2}\times 10^{4} = \qty{1.26e-3}{cm^3}$,
masa de \qty{1.4}{mg} por gramo. Por hectárea ($2.6\times 10^{9}$ g):
\qty{2.6}{t} de bacterias y \qty{3.6}{t} de hongos.
\end{solution}

\begin{solution}{exo:b2:living-soil:9}
$H^{*} = 1.5/0.02 = \qty{75}{t/ha}$ de carbono. Con arado: $1.5/0.04
= \qty{37.5}{t/ha}$; pérdida de \qty{37.5}{t/ha}, cuya mitad desaparece al
cabo de $\ln 2/0.04 = 17$ años.
\end{solution}

\begin{solution}{exo:b2:living-soil:10}
Fosfato: $\sqrt{2\times 10^{-13}\times 8.64\times 10^{6}} =
\qty{1.3}{mm}$ en una temporada --- una raíz a \qty{1}{cm} de la siguiente
alcanza una octava parte del \omterm{def:b2:living-soil:soil}{suelo} que las separa, mientras que unas hifas
separadas \qty{0.1}{mm} lo alcanzan todo. Nitrato: \qty{4}{cm}, más que la
separación entre raíces: llega por sí solo a la raíz (y por \omterm{def:b2:biogeochemical-cycles:reservoir}{flujo} de masa
con el agua), de modo que las micorrizas aportan poco para el nitrato y
mucho para el fosfato.
\end{solution}

\begin{solution}{exo:b2:living-soil:11}
Masa: $0.25\times 10^{4}\times 1.3 = \qty{3250}{t/ha}$. Pérdida neta de
\qty{19}{t/ha} al año: duración de 170 años. Carbono perdido: $20\times
0.03 = \qty{0.6}{t/ha}$ al año. Sin labranza, la pérdida neta es de
\qty{1}{t/ha}: 3250 años.
\end{solution}

\begin{solution}{exo:b2:living-soil:12}
La hojarasca se renueva en uno o dos años, el \omterm{def:b2:living-soil:soil}{humus} en cincuenta, el
horizonte A mineral en miles (con una formación de \qty{0.1}{mm} al año), y
la sal se acumula en décadas y se drena en décadas si hay drenaje. Un
agricultor ve responder la hojarasca y el cultivo en una temporada, y el
\omterm{def:b2:living-soil:soil}{humus} a lo largo de su vida profesional (un descenso de un tercio tarda
treinta años; una recuperación, lo mismo); la \omterm{prop:b2:living-soil:erosion}{erosión del suelo} mineral es
invisible en una vida e irreversible en diez; la sal es visible en una
generación y, sin drenaje, permanente. Las variables lentas del \omterm{def:b2:living-soil:soil}{suelo} son
las que ninguna contabilidad anual registra.
\end{solution}

\begin{solution}{pb:b2:living-soil:1}
\textbf{1.} $0.25\times 10^{4}\times 1.3 = \qty{3250}{t/ha}$;
carbono, $\qty{81}{t/ha}$.
\textbf{2.} $81/12 = \qty{6.8}{t/ha}$ de nitrógeno.
\textbf{3.} Bacterias, \qty{1}{mg/g}, $\qty{3.25}{t/ha}$; hifas,
$\pi(2\times 10^{-4})^{2}\times 2\times 10^{4}\times 1.1 =
\qty{2.8}{mg/g}$, \qty{9}{t/ha}: biomasa microbiana de unas
\qty{12}{t/ha}.
\textbf{4.} Unas \qty{14}{t/ha} de materia viva (masa fresca), alrededor
del \qty{17}{\%} del carbono del \omterm{def:b2:living-soil:soil}{humus} --- un pequeño porcentaje de él si
se cuenta como carbono, puesto que los organismos son sobre todo agua.
\textbf{5.} Carbono microbiano, alrededor de la mitad de \qty{12}{t},
\qty{6}{t}; con una renovación de dos veces al año y un \qty{40}{\%}
retenido, consumen $2\times 6/0.4 = \qty{30}{t}$ de carbono y respiran
\qty{18}{t/ha} de carbono, \qty{66}{t} de $\mathrm{CO_2}$ (una estimación
burda: la retención y la renovación son aproximadas).
\textbf{6.} El triple de la producción primaria neta del cultivo, de
\qty{5}{t}: la estimación es demasiado alta, porque las hipótesis (toda la
biomasa renovándose dos veces, nada de ella en latencia) son generosas; la
respiración real del \omterm{def:b2:living-soil:soil}{suelo} bajo un cultivo es del orden de la propia
producción del cultivo, varias toneladas de carbono por hectárea al año.
El \omterm{def:b2:living-soil:soil}{suelo} respira tanto como las plantas que tiene encima.
\textbf{7.} Carbono: $5\times 0.45 = \qty{2.25}{t}$; nitrógeno: $5\times
0.03 = \qty{150}{kg}$; $\mathrm{C:N} = 15$.
\textbf{8.} Mineralizan: los microorganismos necesitan $2250\times 0.4/8 =
\qty{112}{kg}$ de nitrógeno y los restos contienen 150, de modo que se
liberan netos \qty{38}{kg/ha}.
\textbf{9.} $1 - e^{-0.5} = \qty{39}{\%}$ al cabo de 3 meses; $1 - e^{-2}
= \qty{86}{\%}$ al cabo de un año.
\textbf{10.} $0.39\times 38 = \qty{15}{kg/ha}$.
\textbf{11.} El cultivo necesita \qty{150}{kg}, la mayor parte en sus dos
primeros meses; los restos dan \qty{15}{kg} en tres meses y \qty{38}{kg}
en total --- un complemento, no un suministro. Su valor está también en el
nitrógeno que fijó el trébol y que contienen los propios \qty{150}{kg} de
los restos, la mayor parte del cual pasa al \omterm{def:b2:living-soil:soil}{humus} y vuelve a lo largo de
los años.
\textbf{12.} Paja: carbono, \qty{2250}{kg}; nitrógeno, $2250/80 =
\qty{28}{kg}$; los microorganismos necesitan \qty{112}{kg}, de modo que se
inmovilizan \qty{84}{kg/ha} del \omterm{def:b2:living-soil:soil}{suelo} --- más de la mitad de lo que
necesita el cultivo. El agricultor debe añadir nitrógeno con la paja,
enterrarla mucho antes de la siembra o dejarla en la superficie, donde se
descompone despacio.
\textbf{13.} Con arado: $1.2/0.025 = \qty{48}{t/ha}$; sin labranza:
$1.2/0.015 = \qty{80}{t/ha}$.
\textbf{14.} $H(t) = 80 - 32\,e^{-0.015t}$: ganancia de \qty{4.5}{t/ha}
al cabo de 10 años y de \qty{17}{t/ha} al cabo de 50.
\textbf{15.} $\mathrm{d}H/\mathrm{d}t = 0.015\times 32 =
\qty{0.48}{t/ha}$ de carbono el primer año, \qty{1.8}{t} de
$\mathrm{CO_2}$.
\textbf{16.} $0.48\times 1.5\times 10^{9} = \qty{0.7}{GtC}$ el primer año,
el \qty{7}{\%} de las emisiones fósiles; en el quincuagésimo año, la
velocidad ha bajado a $0.48\,e^{-0.75} = \qty{0.23}{t/ha}$,
\qty{0.34}{GtC}.
\textbf{17.} La ganancia es la aproximación a un nuevo \omterm{def:b2:biogeochemical-cycles:reservoir}{estado
estacionario}: tiende a cero cuando $H$ alcanza \qty{80}{t/ha}, y toda la
ganancia vuelve al aire en unas décadas si el campo se ara de nuevo ---
un sumidero en el \omterm{def:b2:living-soil:soil}{suelo} es una transferencia única y reversible, no una
compensación permanente de unas emisiones que continúan.
\textbf{18.} $k_h = 0.02$: $H^{*} = \qty{60}{t/ha}$, una pérdida de
\qty{20}{t/ha}, mayor que la ganancia de los primeros cincuenta años sin labranza: el calentamiento puede
deshacer lo logrado con el manejo.
\textbf{19.} Con arado: $15 - 0.5 = \qty{14.5}{t/ha}$ al año; sin
labranza: \qty{1}{t/ha}.
\textbf{20.} $3250/14.5 = 220$ años; $3250/1 = 3250$ años.
\textbf{21.} Carbono: $15\times 0.025 = \qty{0.38}{t/ha}$; nitrógeno:
$0.38/12 = \qty{31}{kg/ha}$ al año.
\textbf{22.} La descomposición del \omterm{def:b2:living-soil:soil}{humus} en el \omterm{def:b2:biogeochemical-cycles:reservoir}{estado estacionario} con
arado es $k_hH^{*} =
I = \qty{1.2}{t/ha}$ al año; la erosión añade otro tercio, pero, a
diferencia de la descomposición, no la compensa ningún aporte --- vacía la
reserva.
\textbf{23.} En parte. El carbono erosionado que queda enterrado en un
sedimento o en un embalse puede quedar protegido de la descomposición
durante mucho tiempo (un sumidero por enterramiento); pero los agregados se
rompen durante el transporte, gran parte del carbono se oxida por el
camino, y la fertilidad perdida del campo se sustituye con abonos cuya
fabricación emite carbono. El signo global del efecto de la erosión sobre
el carbono sigue siendo objeto de debate.
\textbf{24.} Hojarasca: $1/k = 0.5$ años; \omterm{def:b2:living-soil:soil}{humus}: $1/k_h = 40$ años con
arado y 67 sin labranza; \omterm{def:b2:living-soil:soil}{suelo} mineral: $3250/0.5 =
6500$ años respecto de la formación, y 220 años respecto de la pérdida con
arado.
\textbf{25.} Reserva de carbono de \qty{81}{t/ha} en el horizonte A;
duración de 220 años con arado y de 3250 sin labranza; los restos de trébol
aportan unos \qty{38}{kg/ha} de nitrógeno mineral, 15 de ellos en los tres
primeros meses.
\end{solution}
